% Draft conclusion for JBES revision, v3 (rewritten 2026-06-07). Bellemare
% formula: summary -> limitations -> implications -> future work. Reframed off the
% dead renter-share moderator (a pre-winsor artifact, now NS and NYC-leverage-
% driven) and the dead Shen-Baur "sign disagreement" (the oriented channels are
% near-collinear, r=0.94). New thesis = framework + a positive, robust-across-
% encodings, heterogeneous, not-cleanly-causal text-price signal. Numbers locked
% against the canonical files and the v3 body sections.

\section{Conclusion}
\label{sec:conclusion}

This paper has asked whether the predictive gains that text embeddings deliver in
automated valuation models reflect causal semantic content or spatial confounding
mediated through language, and whether the answer travels across metropolitan
housing markets. Working with $3{,}851$ deduplicated Redfin listings across twelve
metropolitan areas, a clean panel from which we exclude no market and whose
confounders we winsorize at the median plus or minus five median absolute
deviations to neutralize parcel data-entry artifacts, we built a framework that
joins a partially-linear Double Machine Learning estimator, LEACE concept erasure
with an iterative nonlinear polish, a counterfactual rewriting intervention with
two open-weight language models, and a finite-sample calibration study. The
central empirical finding is that listing text carries a positive
per-standard-deviation association with log sale price that is robust to the
choice of text encoding but heterogeneous across markets. The Doc2Vec uniqueness
channel pools to $+0.072$ and the sentence-BERT principal-component axis to
$+0.156$, and the two are nearly collinear across markets at a Pearson
correlation of $+0.94$, so they recover one signal under two representations
rather than two independent observations of it. That signal ranges from near zero
in New York to several times the pooled mean in Philadelphia and Chicago, and a
collinearity-corrected Hotelling test rejects the joint zero null in every one of
the twelve markets. The counterfactual intervention sharpens the reading. A
contrast that strips stylistic content while holding the model's inferred location
fixed isolates a positive style effect of $+0.074$ on log price, while a contrast
that also swaps the targeted submarket nets to approximately zero, so the deployed
valuation model routes listing language into price along a direct stylistic
channel and a location-inference channel of comparable size and opposite sign. We
searched for a metropolitan-level demographic moderator of the pooled effect and
found none that survives multiplicity adjustment or leave-one-out leverage
diagnostics at twelve markets, and we report that null rather than a moderator
that an earlier draft had reported before a confounder correction reversed its
sign.

\paragraph{Limitations.} Several limitations bound the reading of these results.
The estimates are associational rather than cleanly causal. The
robustness-value sensitivity analysis of Section~\ref{sec:sensitivity-12} shows
that a modest unmeasured confounder, on the order of the spatial covariates we do
observe, would suffice to overturn the per-market estimates, so we do not claim to
have isolated a structural semantic effect, only to have separated the linearly
spatial component of the text signal from the remainder and bounded the strength
of confounding the remainder could still carry. The corpus is restricted to the
twelve metropolitan areas served by the Redfin scraper at the time of collection,
which biases the sample toward larger and more competitive markets and away from
the rural and small-metro regimes in which language-mediated spatial confounding
is plausibly different in character. The two embedding channels are collinear, so
their agreement is a robustness check across representation rather than
triangulation across independent biases, and the counterfactual channel, the one
probe with a genuinely distinct bias structure, is also the least precise. Finally
the panel is small enough at twelve markets that one high-leverage market, New
York, dominates the moderator diagnostics, so the null moderator result should be
read as the absence of a moderator detectable at this sample size rather than as a
demonstrated absence of moderation.

\paragraph{Implications.} The clearest implication is methodological. Predictive
fit cannot distinguish a text embedding that prices genuine semantic content from
one that launders neighborhood geography through language, because the two are
observationally identical in a predictive metric, and our results show that a
large share of the apparent text signal is exactly the spatial information that
linear concept erasure removes. An analyst or regulator who treats a text-AVM's
predictive gain as evidence of causal semantic valuation is therefore
overstating what the gain establishes, and the deconfounding step we describe,
erasing the location concept and re-estimating, is a low-cost diagnostic that
separates the two before either is trusted. Because the residual signal that
survives deconfounding is heterogeneous across markets and sensitive to unmeasured
confounding, the appropriate posture toward NLP-augmented valuation is neither to
dismiss the text channel nor to treat it as a settled causal input, but to audit
it market by market with the design-based tools this paper assembles.

\paragraph{Future research.} Several extensions follow. A panel of substantially
more than twelve markets would give the moderator analysis the leverage it lacks
here and would let the cross-market heterogeneity we document be explained rather
than only described. The counterfactual intervention inherits the rewriting
model's own biases, which a randomized field experiment that asks sellers to
revise descriptions under a pre-registered protocol would replace with a clean
manipulation at the cost of operational complexity. Real recorded transaction
prices in place of the scraped sale prices would tighten the outcome and remove
any residual concern about price measurement, and pairing the causal-residual
detector this paper isolates with the demographic-stratification machinery of the
fairness-audit literature would let practitioners locate where any
language-mediated mispricing that survives deconfounding is large enough to matter
for borrowers.
